\clearpage
\section{일반 방정식과 전체 대칭군}
\label{sec:generic-equation}

\begin{learninggoal}
변수 자체를 근으로 갖는 일반 방정식을 구성하고 그 갈루아군이 $S_n$임을 증명한다. 대칭 유리함수의 고정체가 기본대칭식으로 생성됨을 이해한다.
\end{learninggoal}

낮은 차수의 특별한 다항식에서는 갈루아군이 작은 부분군으로 줄어들 수 있다. 그렇다면 계수 사이에 아무 특별한 관계도 없는 ``일반적인'' $n$차방정식의 갈루아군은 무엇일까? 답은 전체 대칭군 $S_n$이다.

\subsection{근을 독립변수로 두기}

체 $k$와 대수적으로 독립인 변수
\[
  T_1,\dots,T_n
\]
을 잡고
\[
  L=k(T_1,\dots,T_n)
\]
이라 하자. 각 $\pi\in S_n$은 변수들을 순열하여
\[
  \pi\bigl(R(T_1,\dots,T_n)\bigr)
  =R(T_{\pi(1)},\dots,T_{\pi(n)})
\]
인 $k$-자동동형을 정의한다. 이 작용의 고정체를
\[
  K=L^{S_n}
\]
라 하자. 아르틴의 고정체 정리에 의해
\[
  L/K
\]
는 갈루아 확장이고
\[
  \Gal(L/K)=S_n,
  \qquad
  [L:K]=n!.
\]

기본대칭식들을
\[
\begin{aligned}
  s_1&=T_1+\cdots+T_n,\\
  s_2&=\sum_{i<j}T_iT_j,\\
  &\ \vdots\\
  s_n&=T_1\cdots T_n
\end{aligned}
\]
라 하자. 그러면
\[
  F(X)=\prod_{i=1}^n(X-T_i)
  =X^n-s_1X^{n-1}+s_2X^{n-2}-\cdots+(-1)^ns_n.
\]
모든 $s_i$는 변수의 순열에 불변이므로
\[
  k(s_1,\dots,s_n)\subseteq K.
\]

\begin{theorem}[대칭 유리함수의 고정체]
\[
  k(T_1,\dots,T_n)^{S_n}
  =k(s_1,\dots,s_n).
\]
또한 $s_1,\dots,s_n$은 $k$ 위에서 대수적으로 독립이다.
\end{theorem}

\begin{proof}
$F$는 $k(s_1,\dots,s_n)$ 위에서 $T_1,\dots,T_n$을 모두 근으로 갖고, $L$은 그 분해체이다. 근 순열 표현에 의해
\[
  [L:k(s_1,\dots,s_n)]\le n!.
\]
한편
\[
  k(s_1,\dots,s_n)\subseteq K
\]
이고 $[L:K]=n!$이므로 탑 법칙에서 두 중간체는 같아야 한다.

대수적 독립성을 보이기 위해 새로운 독립변수 $S_1,\dots,S_n$을 잡고 일반 다항식
\[
  X^n-S_1X^{n-1}+\cdots+(-1)^nS_n
\]
의 분해체에서 근들을 $t_1,\dots,t_n$이라 하자. 비에타 관계로 $S_i$는 $t_j$들의 기본대칭식이다. 만약 $s_i$ 사이에 영이 아닌 다항식 관계가 있다면 같은 관계를 $S_i$에 대입해 얻지만, $S_i$들은 독립변수이므로 모순이다.
\end{proof}

\begin{corollary}[일반 $n$차방정식의 갈루아군]
독립변수 $S_1,\dots,S_n$에 대해
\[
  P(X)=X^n+S_1X^{n-1}+\cdots+S_n
  \in k(S_1,\dots,S_n)[X]
\]
라 하자. 그러면 $P$는 분리 기약다항식이고
\[
  \Gal(P/k(S_1,\dots,S_n))\cong S_n.
\]
\end{corollary}

\begin{proof}
부호를 조정한 기본대칭식과 $S_i$를 대응시키면 앞의 구성과 동형인 분해체를 얻는다. 그 갈루아군은 $S_n$이다. $S_n$이 근들 위에서 추이적으로 작용하므로 $P$는 기약이고, 독립변수인 근들이 서로 다르므로 분리이다.
\end{proof}

\begin{remark}[``일반적''이라는 말의 의미]
이 정리는 모든 수치계수 $n$차다항식의 갈루아군이 $S_n$이라는 뜻이 아니다. 계수들을 독립변수로 놓아 어떤 특별한 대수적 관계도 강제하지 않았을 때 전체 대칭군이 나타난다는 뜻이다. 특정 계수로 전문화하면 판별식이 제곱이 되거나 분해식이 인수분해되어 갈루아군이 작은 부분군으로 내려갈 수 있다.
\end{remark}

\subsection{대칭다항식의 첫 정리}

\begin{theorem}[기본대칭식 정리: 체 계수 버전]
대칭다항식
\[
  H\in k[T_1,\dots,T_n]
\]
은 유일하게
\[
  H=G(s_1,\dots,s_n)
\]
꼴로 쓸 수 있다. 여기서 $G\in k[Y_1,\dots,Y_n]$이다.
\end{theorem}

\begin{proofidea}
사전식 순서에서 가장 큰 단항식을 선택한다. 대칭성 때문에 그 지수는 감소하는 순서로 배열되며, 적당한
\[
  s_1^{m_1-m_2}s_2^{m_2-m_3}\cdots s_n^{m_n}
\]
을 빼면 최고항이 제거된다. 이 과정을 유한번 반복한다. 유일성은 $s_i$의 대수적 독립성에서 따른다.
\end{proofidea}

이 정리는 근들에 대칭인 표현을 계수로 바꾸는 원리이다. 다음 장에서 판별식
\[
  \prod_{i<j}(\alpha_i-\alpha_j)^2
\]
이 계수들의 다항식으로 표현된다는 사실을 정당화하는 데 사용한다.

\begin{example}
두 변수에서
\[
  T_1^2+T_2^2
  =(T_1+T_2)^2-2T_1T_2
  =s_1^2-2s_2.
\]
세 변수에서는
\[
  T_1^2+T_2^2+T_3^2=s_1^2-2s_2.
\]
근들의 제곱합이 계수로 표현되는 가장 간단한 예이다.
\end{example}

\begin{keyidea}
전체 대칭군은 ``근의 이름만 바뀌고 계수로 표현되는 정보는 변하지 않는'' 일반 상황의 대칭이다. 기본대칭식은 근의 세계와 계수의 세계를 연결하는 좌표계이다.
\end{keyidea}

\begin{checkpoint}
\begin{enumerate}[label=(\alph*)]
  \item $S_n$이 $k(T_1,\dots,T_n)$에 충실하게 작용하는 이유를 설명하라.
  \item $[L:k(s_1,\dots,s_n)]\le n!$은 어떤 분해체 상계에서 나오는가?
  \item $T_1^3+T_2^3$을 $s_1=T_1+T_2$, $s_2=T_1T_2$로 표현하라.
\end{enumerate}
\end{checkpoint}
